\scanpage{79}
% archive foliation: 78; typescript printed page number follows
\begin{center}
- 16 -
\end{center}

ce qui prouve que la série considérée est absolument convergente pour
$\|u\|<1/5a$. Le fait que $(h(u))^2=e+u$ résulte de ce que la formule
$(1-\sum_{n=1}^{\infty}\binom{2n}{n}x^n)^2=1-4x$ est vraie lorsque
$E=K=\mathbb R$ (\uline{Fonct.var.réelle}, chap.III,§2)
(cf.n°2, remarque suivant la prop.2).

Reste à montrer que la fonction $g$ est la seule fonction continue et égale
à $e$ à l'origine, dont le carré soit $f$. Or, la relation $g^2-g_1^2=0$
s'écrit $(g-g_1)(g+g_1)=0$ puisque $E$ est commutative ; mais comme
$g(0)+g_1(0)=2e$, $g+g_1$ est inversible dans un voisinage de l'origine,
donc on a $g-g_1=0$ dans ce voisinage, ce qui achève la démonstration.

\begin{center}
§ 2. La méthode des majorantes.
\end{center}

\uline{1. Majorantes d'une série entière}.

DEFINITION 1.- \uline{Etant donnée une série entière}
$(\unsure{\mathbf c|C}_{n_1n_2\ldots n_p}z_1^{n_1}\ldots z_p^{n_p})$
\uline{à coefficients dans} $E$, \uline{par rapport aux variables}
$z_i\in K$ $(1\leq i\leq p)$,
\uline{on dit qu'une série entière}
$(C_{n_1n_2\ldots n_p}x_1^{n_1}\ldots x_p^{n_p})$
\uline{à coefficients \unsure{positifs}}
\uline{par rapport aux variables réelles} $x_i$ $(1\leq i\leq p)$,
\uline{est une majorante de \unsure{la} série entière donnée, si pour tout système d'exposants}
$(n_1,n_2,\ldots,n_p)$, \uline{on a}
$\|\unsure{\mathbf c|C}_{n_1n_2\ldots n_p}\|\leq C_{n_1n_2\ldots n_p}$.

Il est clair que la somme (resp. le produit) de majorantes de deux séries
entières est une majorante de la somme (resp. du produit) de ces deux séries.

Si une majorante d'une série entière est sommable pour
$x_1=r_1>0,\ldots,x_p=r_p>0$, il est clair que la série entière donnée
est absolument sommable dans le polycylindre fermé défini par
$|z_i|\leq r_i$ $(1\leq i\leq q)$.
Inversement, supposons qu'une série entière
$(\unsure{\mathbf c|C}_{n_1\ldots n_p}z_1^{n_1}\ldots z_p^{n_p})$
soit absolument sommable dans le polycylindre fermé
